% ===== PRÉAMBULE CANONIQUE — à coller VERBATIM en tête de CHAQUE .tex =====
\documentclass[11pt,a4paper]{article}
\usepackage[utf8]{inputenc}\usepackage[T1]{fontenc}\usepackage{lmodern}
\usepackage[french]{babel}
\usepackage{amsmath,amssymb,mathtools}\usepackage{icomma}
\usepackage[version=4]{mhchem}          % \ce{...} : équations & formules
\usepackage{siunitx}\sisetup{locale=FR} % \qty{}{}, \si{}, \ang{}
\usepackage{chemfig}                    % \chemfig{...} : structures développées
\usepackage{xcolor}\usepackage{colortbl}
\usepackage{tikz}\usepackage{pgfplots}\pgfplotsset{compat=1.18}
\usepackage[most]{tcolorbox}\usepackage{enumitem}\usepackage{array,booktabs}
\usepackage[a4paper,margin=2.2cm]{geometry}
\usetikzlibrary{arrows.meta,calc,angles,quotes,positioning,patterns,backgrounds,shapes.geometric}
\definecolor{primary}{RGB}{222,49,99}\definecolor{primarydark}{RGB}{150,22,55}
\definecolor{defcol}{RGB}{0,114,178}\definecolor{methcol}{RGB}{52,92,148}
\definecolor{excol}{RGB}{0,135,90}\definecolor{attncol}{RGB}{200,40,55}
\tcbset{boxbase/.style={enhanced,breakable,boxrule=0.9pt,arc=2.5pt,
  left=3mm,right=3mm,top=2.3mm,bottom=2.3mm,fonttitle=\bfseries,coltitle=white}}
% --- boîtes TUTORIEL (noms EXACTS, ne pas renommer) ---
\newtcolorbox{cle}[1][]{boxbase,colback=defcol!6,colframe=defcol,
  title={\textsc{À retenir}\ifx\relax#1\relax\relax\else\textmd{ : \itshape #1}\fi}}
\newtcolorbox{methode}[1][]{boxbase,colback=methcol!7,colframe=methcol,sharp corners,
  title={\textsc{Méthode}\ifx\relax#1\relax\relax\else\textmd{ --- \itshape #1}\fi}}
\newtcolorbox{exemplebox}[1][]{boxbase,colback=excol!5,colframe=excol,
  title={\textsc{Exemple}\ifx\relax#1\relax\relax\else\textmd{ : \itshape #1}\fi}}
\newtcolorbox{piege}[1][]{boxbase,colback=attncol!6,colframe=attncol,
  title={\textsc{Piège}\ifx\relax#1\relax\relax\else\textmd{ : \itshape #1}\fi}}
\newtcolorbox{remarque}[1][]{boxbase,colback=black!4,colframe=black!45,
  title={\textsc{Remarque}\ifx\relax#1\relax\relax\else\textmd{ : \itshape #1}\fi}}
% --- boîtes EXERCICE / CORRIGÉ (noms EXACTS, auto-comptées) ---
\newtcolorbox[auto counter]{exo}[1][]{boxbase,colback=primary!4,colframe=primary,
  title={\textsc{Exercice}~\thetcbcounter\ifx\relax#1\relax\relax\else\textmd{ --- \itshape #1}\fi}}
\newtcolorbox[auto counter]{corrige}[1][]{boxbase,colback=excol!4,colframe=excol,
  title={\textsc{Corrigé}~\thetcbcounter\ifx\relax#1\relax\relax\else\textmd{ --- \itshape #1}\fi}}
\newcommand{\R}{\mathbb{R}}\newcommand{\N}{\mathbb{N}}\newcommand{\Z}{\mathbb{Z}}\newcommand{\ds}{\displaystyle}
\begin{document}
\begin{corrige}[Nommer entièrement une molécule à deux fonctions]
\begin{enumerate}[label=\alph*)]
  \item Un $-\ce{COOH}$ (\textbf{acide carboxylique}) et un $-\ce{OH}$ (\textbf{alcool secondaire}). Priorité : $-\ce{COOH} > -\ce{OH}$.
  \item L'acide donne le suffixe (\textit{acide \dots-oïque}). L'alcool devient le préfixe \textbf{hydroxy-}.
  \item Chaîne de $4$~carbones (butanoïque), le carbone du $-\ce{COOH}$ est le n$^\circ$1 : \ce{C1OOH-C2H2-C3H(OH)-C4H3}. Le $-\ce{OH}$ est en position 3. Nom : \textbf{acide 3-hydroxybutanoïque}.
\end{enumerate}
\end{corrige}

\begin{corrige}[Trois fonctions, une seule principale]
\begin{enumerate}[label=\alph*)]
  \item \textbf{alcène} ($\ce{C=C}$), \textbf{alcool} ($-\ce{OH}$) et \textbf{cétone} ($-\ce{CO}-$).
  \item Priorité décroissante : cétone $-\ce{CO}- \;>\; -\ce{OH} \;>\; \ce{C=C}$.
  \item La \textbf{cétone} fixe le suffixe \textit{-one}. L'alcool devient le préfixe \textbf{hydroxy-} ; la double liaison est signalée par l'infixe \textbf{-én-} dans le nom de la chaîne.
  \item En numérotant depuis le $\ce{CH3}$ pour donner à la cétone le plus petit indice : cétone en 2, $-\ce{OH}$ en 3, double liaison en 4. Nom : \textbf{3-hydroxypent-4-én-2-one}.
\end{enumerate}
\end{corrige}

\begin{corrige}[Du nom à la structure : le préfixe « oxo- »]
\begin{enumerate}[label=\alph*)]
  \item «\,butano\,» $\Rightarrow$ chaîne de $4$~carbones ; «\,acide \dots-oïque\,» $\Rightarrow$ fonction principale \textbf{acide carboxylique} (carbone n$^\circ$1).
  \item «\,3-oxo\,» signale un carbonyle $\ce{C=O}$ (ici une \textbf{cétone}) sur le carbone n$^\circ$3, en fonction secondaire.
  \item Formule semi-développée : \ce{CH3-CO-CH2-COOH}. Formule brute : \ce{C4H6O3}.
\end{enumerate}
\end{corrige}

\begin{corrige}[Acide et aldéhyde en compétition]
\begin{enumerate}[label=\alph*)]
  \item un $-\ce{CHO}$ (\textbf{aldéhyde}) et un $-\ce{COOH}$ (\textbf{acide carboxylique}).
  \item L'\textbf{acide carboxylique} est principal : dans l'échelle, $-\ce{COOH} > -\ce{CHO}$.
  \item L'aldéhyde secondaire apparaît sous le préfixe \textbf{oxo-}. La chaîne compte $3$~carbones, le $-\ce{COOH}$ est en n$^\circ$1, le $-\ce{CHO}$ en n$^\circ$3 : nom \textbf{acide 3-oxopropanoïque}.
  \item Formule brute : \ce{C3H4O3}.
\end{enumerate}
\end{corrige}

\end{document}
